% =============================================================================
\subsection{对平台治理与经验校准的启示}
\label{sec:governance}
上述结果对信息生态系统的治理提出三条具体的、且在某种程度上反直觉的启示。

\paragraph{(i) 同步性是风险指标，而不是健康指标。}
平台与观察者通常把高度一致——趋势的整齐划一、集体注意力的快速同步转向——
解读为公共领域运行良好的证据。图~\ref{fig:vuln} 的脆弱性窗口颠倒了这一
读法：使群体相干响应最大化的那个噪声水平，\emph{同时}也是错误信念被最有效
组织起来的水平。一个实际的后果是：集体相干性的度量（总体情绪的波动幅度、
分享行为的同步性）可以被作为易感性的\emph{预警指标}加以监测，正如临界慢化
（critical slowing down）被用作生态与气候临界点附近的预警信号一样。这提供
了一条从"事后辟谣"走向"事前预警"的路径。

\paragraph{(ii) 在系统最可激发之处进行干预。}
辟谣热力图（图~\ref{fig:heat}）表明，真相偏置的边际效应在脆弱性窗口内最大。
由于那里的信念最"可组织"，同样的辟谣投入能买到比在沉寂区或混沌区更大的
流行率下降。这支持\emph{自适应}地分配审核与事实核查资源——在集体可激发性高的
时期集中投入——而不是维持一个恒定、无针对性的预算。这是本文最具操作性的
政策含义：\emph{最危险的信息环境同时也是最可治理的信息环境}。

\paragraph{(iii) 靶向波前，而不只是源头。}
图~\ref{fig:mis-wf} 的空间波前显示错误信念是作为一个从焦点种子发出的\emph{
相干波前}扩散的，而不是作为大量独立的皈依事件。因此，仅移除首发账号或
"对源头做事实核查"的干预只处理了动力学中趋于零的一部分；遏制一场共振级联
需要作用于波前正在推进\emph{进去}的那个邻域。这与流行病学中的环形接种
（ring vaccination）思路同构：在传播前沿而非起点处投入资源。

\paragraph{通往经验校准的路径。}
本模型是一个刻意的机制性理想化，但其参数具有明确的经验对应物：稳定指数
$\alpha$ 可由级联规模或事件间隔分布的尾指数估计
\citep{barabasi2005,goel2012}；噪声强度 $D$ 由总体活动量估计；耦合
$J_0,\sigma$ 由转发/提及网络所揭示的交互范围估计；真相偏置 $C$ 由纠正信息
相对于原始主张的可测量触达估计。从平台数据中估计这些参数——而非假定它们——
将把本框架从思想实验转变为经过校准的工具，并使预言
$D^{*}\approx\Delta U/2$ 能够被拿来与观测到的易感性升高期相对照。我们认为
这是最有价值的下一步。

